% ===================================================================
%  Αρχείο: 37101_SOLUTION.tex (Fragment)
%  Αυτόματη μετατροπή μέσω doc2latex.py
% ===================================================================

ΘΕΜΑ 4

Δίνεται τρίγωνο ΑΒΓ με γωνίες Β και Γ οξείες και Δ, Μ και Ε τα μέσα των πλευρών του ΑΒ, ΑΓ και ΒΓ αντίστοιχα. Στις μεσοκάθετες των ΑΒ και ΒΓ και εκτός του τριγώνου ΑΒΓ θεωρούμε σημεία Ζ και Η αντίστοιχα, τέτοια ώστε $\Delta Ζ = \ \frac{ΑΒ}{2}$ και $ΕΗ = \ \frac{Β\Gamma}{2}$.

\begin{alist}
\item Να αποδείξετε ότι:

.}
\end{alist}
\begin{alist}
\item
  Το τετράπλευρο ΒΔΜΕ είναι παραλληλόγραμμο. \monades{5}
\item
  Τα τρίγωνα ΖΔΜ και ΕΜΗ είναι ίσα. \monades{10}

\item Αν τα σημεία Ζ, Δ, Ε είναι συνευθειακά, να αποδείξετε ότι η γωνία $\widehat{Α} = \ 90^{ο}$. \monades{10}

\includesvg[width=5.32292in,height=3.94375in]{/home/spyros/Μαθηματικά/Trapeza_Thematwn/A_Lykeiou/Geometry/extracted/37101_SOLUTION/media/image2.svg}
\end{alist}
